% Source-derived chapter 01: Introduction
% Original source: pw-conjecture-gl-n
\chapter{Introduction}
\label{pw-conjecture-gl-n-chapter-01}
\source{pw-conjecture-gl-n}

\begin{remark}[Source section]
\label{pw-conjecture-gl-n-chapter-01-source}
\source{pw-conjecture-gl-n}
\source{pw-conjecture-gl-n}
This chapter reproduces the corresponding section of the retrieved
LaTeX source, including its definitions, statements, examples, and
proofs. It has not yet been translated into Lean.
\notready
\end{remark}

% The source section title was: Introduction
Throughout, we work over the complex numbers $\BC$. 

\subsection{The P=W conjecture}
The purpose of this paper is to present a proof of the $P=W$ conjecture by de Cataldo--Hausel--Migliorini \cite{dCHM1} for arbitrary rank $n$ and genus $g \geq 2$.

Let $C$ be a nonsingular irreducible projective curve of genus $g \geq 2$. For two coprime integers $n \in \BZ_{\geq 1}$ and $d \in \BZ$, there are two moduli spaces $M_{\mathrm{Dol}}$ and $M_B$, called the Dolbeault and the Betti moduli spaces, attached to $C,n,$ and $d$. 


The Dolbeault moduli space parameterizes stable Higgs bundles $(\CE, \theta)$ on $C$, where $\CE$ is a vector bundle on $C$ of rank $n$ and degree $d$, $\theta: \CE \to \CE \otimes \Omega_C^1$ is a Higgs field, and stability is defined with respect to the slope $\mu(\CE,\theta) ={\mathrm{deg}(\CE)}/{\mathrm{rk(\CE)}}$. The moduli space $M_{\mathrm{Dol}}$ admits the structure of a completely integrable system
\[
h: M_{\mathrm{Dol}} \to A := \bigoplus_{i=1}^n H^0(C, {\Omega_C^{1}}^{\otimes i}), \quad (\CE, \theta) \mapsto \mathrm{char.polynomial}(\theta),
\]
which is referred to as the Hitchin system \cite{Hit, Hit1}. The Hitchin map $h$ is surjective and proper; it is also Lagrangian with respect to the canonical holomorphic symplectic form on $M_{\mathrm{Dol}}$ induced by the hyper-K\"ahler metric. The \emph{perverse filtration} is an increasing filtration 
\[
P_0H^*(M_{\mathrm{Dol}}, \BQ)\subset P_1H^*(M_{\mathrm{Dol}}, \BQ) \subset \cdots \subset H^*(M_{\mathrm{Dol}}, \BQ)
 \]
on the (singular) cohomology of $M_{\mathrm{Dol}}$ governed by the topology of the Hitchin system $h$; see Section \ref{sec1.1} for a brief review.

The Betti moduli space $M_B$ is the (twisted) character variety associated with $\mathrm{GL}_n(\BC)$ and degree $d$.  It parametrizes isomorphism classes of irreducible local systems
\[
\rho: \pi_1(C\backslash \{p\}) \rightarrow \mathrm{GL}_n(\BC)
\]
where $\rho$ sends a loop around a chosen point $p$ to 
$e^{\frac{2\pi \sqrt{-1} d}{n}}\mathrm{Id}_n$.  Concretely, we have
\begin{equation*}
M_B := \Big{\{}a_k, b_k \in \mathrm{GL}_n(\BC),~k=1,2,\dots,g: ~~\prod_{j=1}^g [a_j, b_j] = e^{\frac{2\pi \sqrt{-1} d}{n}}\mathrm{Id}_n \Big{\}}\git\mathrm{GL}_n(\BC).
\end{equation*}
It is an affine variety whose mixed Hodge structure admits a nontrivial weight filtration 
\[
W_0 H^*(M_B, \BQ) \subset W_1H^*(M_B, \BQ) \subset \cdots \subset H^*(M_B, \BQ).
\]

Non-abelian Hodge theory \cite{Simp, Si1994II} gives a diffeomorphism between the two very different algebraic varieties $M_{\mathrm{Dol}}$ and $M_B$, which canonically identifies their cohomology:
\begin{equation}\label{NAH}
\source{pw-conjecture-gl-n}
 H^*(M_{\mathrm{Dol}}, \BQ)  = H^*(M_{B}, \BQ).
\end{equation}
The $P=W$ conjecture by de Cataldo--Hausel--Migliorini \cite{dCHM1} 
refines the identification (\ref{NAH}); it predicts that the perverse filtration associated with the Hitchin system is matched with the (double-indexed) weight filtration associated with the Betti moduli space. This establishes a surprising connection between topology of Hitchin systems and Hodge theory of character varieties.

\begin{conj}
\source{pw-conjecture-gl-n}
\notready\label{conj}[The P=W conjecture for $\mathrm{GL}_n$, \cite{dCHM1}] For any $k,m\in \BZ_{\geq 0}$, we have
\source{pw-conjecture-gl-n}
\[
P_kH^m(M_{\mathrm{Dol}}, \BQ) = W_{2k}H^m(M_{B}, \BQ) = W_{2k+1}H^m(M_B, \BQ).
\]
\end{conj}

Conjecture \ref{conj} has previously been proven for $n=2$ and arbitrary genus $g\geq 2$ by de Cataldo--Hausel--Migliorini \cite{dCHM1}, and recently for arbitrary rank $n$ and genus 2 by de Cataldo--Maulik--Shen \cite{dCMS}. The compatibility between the $P=W$ conjecture and Galois conjugation on the Betti side was proven in \cite{dCMSZ}, which implies that $P=W$ does not depend on the degree $d$ as long as it is coprime to $n$. We refer to \cite[Section 1]{dCHM1} and the paragraphs following \cite[Theorem 0.2]{dCMS} for discussions concerning connections between Conjecture \ref{conj} and other directions. In particular, by \cite{CDP} and \cite[Section 9.3]{MT}, Conjecture \ref{conj} implies the correspondence between Gopakumar--Vafa invariants and Pandharipande--Thomas invariants \cite[Conjecture 3.13]{MT} for the local Calabi--Yau 3-fold $T^*C \times \BC$ in any curve class $n[C]$. 


The main result of this paper is a full proof of Conjecture \ref{conj}.

\begin{thm}
\source{pw-conjecture-gl-n}
\notready\label{thm}
\source{pw-conjecture-gl-n}
    Conjecture \ref{conj} holds.
\end{thm}


Conjecture \ref{conj} has several variants which, to the best of our knowledge, are still open. These include the version formulated for possibly singular moduli spaces, intersection cohomology, and general reductive groups \cite{dCHM1, dCM, FM}, and the version formulated for moduli stacks \cite{Davison}.  There also exist parabolic versions of Conjecture \ref{conj} and we expect that our argument applies to these settings using parabolic variants of the ingredients here \cite{Mellit, OY}. We refer to \cite{DMS, HLSY, Geo_P=W, Zhang} and references therein for the $P=W$ phenomenon in other settings.


For the case of type $A_{n-1}$ ($\mathrm{GL}_n, \mathrm{PGL}_n, \mathrm{SL}_n$) and a degree $d$ coprime to $n$, Conjecture \ref{conj} (\emph{i.e.} the $P=W$ conjecture for $\mathrm{GL_n}$) is equivalent to the $P=W$ conjecture for $\mathrm{PGL}_n$; see the paragraph after \cite[Theorem 0.2]{dCMS}. The case of $\mathrm{SL}_n$ is more subtle --- it is closely related to the endoscopic decomposition of the $\mathrm{SL}_n$ Hitchin moduli spaces. A systematic discussion on this aspect can be found in \cite[Section 5]{MS}. In the special case when $n=p$ a prime number, \cite[Theorem 0.2]{SLn} implies that the $P=W$ conjecture for $\mathrm{GL}_p$, $\mathrm{SL}_p$, and $\mathrm{PGL}_p$ are all equivalent. In particular, we have the following immediate consequence of Theorem \ref{thm}, which generalizes the result of \cite{dCHM1} for $\mathrm{SL}_2$.

\begin{thm}
\source{pw-conjecture-gl-n}
\notready
The $P=W$ conjecture holds for a curve $C$ of any genus at least two and $\mathrm{SL}_p$ with $p$ a prime number.
\end{thm}






\subsection{Idea of the proof}
Our proof of Conjecture \ref{conj} has 4 major steps.

\begin{enumerate}
    \item[Step 1.] \emph{Strong perversity of Chern classes}: Using work of Markman \cite{Markman}, Shende \cite{Shende}, and Mellit \cite{Mellit}, we may reduce the $P=W$ conjecture to a statement about the interaction between Chern classes of the universal family and the perverse filtration. This allows us to reduce the $P=W$ conjecture to a sheaf-theoretic statement which concerns strong perversity of Chern classes.
    \item[Step 2.] \emph{Vanishing cycle techniques}: We apply the formalism of vanishing cycles to reduce the sheaf-theoretic formulation of Step 1 for the Hitchin system, to the case of twisted Hitchin systems associated with meromorphic Higgs bundles. Our motivation is from the key observation by Ng\^o \cite{Ngo} and Chaudouard-Laumon \cite{CL}, that the decomposition theorem for such twisted Hitchin systems is more manageable. 
    
    Steps 1 and 2 are carried out in Sections 1 and 2.  

    \item[Step 3.]\emph{Global Springer theory}: The global Springer theory of Yun \cite{Yun1, Yun2,Yun3} produces rich symmetries for certain Hitchin moduli spaces; this proves the sheaf-theoretic formulation of Step 1 over the \emph{elliptic locus}. For our purpose, we need a version of global Springer theory for the stable locus over the total Hitchin base. This is described in Section 3.

    \item[Step 4.] \emph{A support theorem}: Lastly, in order to extend part of Yun's symmetries induced by Chern classes from the elliptic locus to the total Hitchin base, we prove a support theorem parallel to \cite{CL} for certain parabolic Hitchin moduli spaces. This is completed in Section 4.
\end{enumerate}

The idea of combining vanishing cycle functors and support theorems was also applied in \cite{MS} to give a proof of the topological mirror symmetry conjecture of Hausel--Thaddeus \cite{HT, GWZ}.

\subsection{Acknowledgements}
We would like to thank Mark Andrea de Cataldo, Bhargav Bhatt, Ben Davison, Jochen Heinloth, and Max Lieblich for various discussions.  We are especially grateful to Zhiwei Yun for explaining his thesis to us, both in his office and at the playground.  We also thank the anonymous referee for careful
reading and useful suggestions on the exposition. J.S. was supported by the NSF grants DMS-2134315 and DMS-2301474.
